\documentclass[10pt,a4paper]{amsart}
\usepackage[margin=19mm]{geometry}
\usepackage{amsmath,amssymb,mathtools}
\usepackage[scr=rsfs,cal=euler]{mathalfa}
\usepackage{xfrac}
\usepackage[unicode,hidelinks,bookmarks=false]{hyperref}
\newcommand{\ie}{i.e.\ }
\newcommand{\cf}{cf.\ }
\newcommand{\resp}{respectively\ }
\newcommand{\Subsection}{\subsection}
\providecommand{\nobreakdash}{\nobreak\hbox{-}}
\setlength{\emergencystretch}{2em}
\multlinegap0pt
\usepackage{amssymb}
\usepackage{xcolor}
\usepackage{braket}
\usepackage{cancel}
\usepackage[normalem]{ulem}
\usepackage[shortlabels]{enumitem}
\newtheorem{proposition}{Proposition}
\DeclareMathOperator{\id}{id}
\title{Multi-entropy in random tensor networks}
\author{Miao Hu, Simon Lin, Ion Nechita}
\date{Source: arXiv:2606.04470v2 (CC BY 4.0)}
\begin{document}
\maketitle

\noindent\emph{Source and licence.} Multi-entropy in random tensor networks. Version arXiv:2606.04470v2 (CC BY 4.0), distributed by arXiv under the Creative Commons Attribution 4.0 International licence, http://creativecommons.org/licenses/by/4.0/, as recorded in the arXiv licence metadata for this exact version and shown on the abstract page https://arxiv.org/abs/2606.04470v2. Full licence notice: \texttt{licenses/arxiv-2606.04470-CC-BY-4.0.txt}.\par\smallskip

\noindent\textit{Editorial scope.} Counted unit: the article's proposition prop:geodesic (source0.tex lines 1870--1884). The article's complete original proof of it is delivered with it (source0.tex lines 1885--1950, one \texttt{proof} environment of 66 lines). No mathematics is written, repaired or re-derived: the statement and the proof are the article's own lines, in the article's own notation, and no retained line is substituted or re-flowed.  No further source-local context is needed: every cross-reference of every retained line resolves inside the two delivered spans.  Nothing was omitted from any retained span, so the package carries no omission record.  No retained line contains a citation, so the wrapper carries no bibliography and no bibliography entry.  No retained line contains a figure environment, an image-inclusion command or drawing code, so no image is delivered and the package declares no asset dependency.  Editorial formatting only, in addition: an AMS \texttt{amsart} preamble that restates the article's own declarations and macros used by the retained lines, namely the environments \texttt{proposition} on separate counters, together with the standard packages those lines need.  The retained environments print in this document as Proposition 1 in order of appearance, whereas the article prints the same items with the numbering of its own full document.  The complete native source of the article as distributed in the arXiv e-print for this exact version is bundled unchanged as \texttt{sources/202049/source0.tex}.
\begin{proposition}
\label{prop:geodesic}
For distinct unordered pairs of terminals $\{i,j\}$ and $\{k,\ell\}$, one has
\begin{equation}
[g_i,g_j]\cap[g_k,g_\ell]
=
\begin{cases}
\{g_m\}, & \{i,j\}\cap\{k,\ell\}=\{m\},\\
\varnothing, & \{i,j\}\cap\{k,\ell\}=\varnothing.
\end{cases}
\label{eq:pairwise-geodesic-interval-intersection}
\end{equation}
Consequently the relative interiors
$[g_i,g_j]\setminus\{g_i,g_j\}$ are pairwise disjoint.
\end{proposition}

\begin{proof}
  Fix $a\ne b$.  Recall that, for $v\in X$, $\tau_v$ denotes the
  translation permutation $\tau_v(x)=x+v$.  By left-invariance of the
  Cayley metric, if
  $\alpha\in[g_a,g_b]$, then
  \begin{equation}
    \beta:=g_a^{-1}\alpha\in[\id,g_a^{-1}g_b]
    =
    [\id,\tau_{e_b-e_a}].
  \end{equation}
  The cycles of $\tau_{e_b-e_a}$ are exactly the affine lines
  \begin{equation}
    x+\langle e_b-e_a\rangle,\qquad x\in X.
  \end{equation}
  The geodesic fact above therefore implies that $\beta$ preserves each
  such affine line.  Equivalently, for every $x\in X$,
  \begin{equation}
    \label{eq:geodesic-affine-line-constraint}
    \alpha(x)-x
    \in e_a+\langle e_b-e_a\rangle
    =
    \{(1-t)e_a+te_b:t\in\mathbb Z_n\}.
  \end{equation}

  Now suppose first that the two unordered pairs share exactly one
  terminal.  Relabeling the three indices, it is enough to consider
  $\{i,j\}$ and $\{i,k\}$, with $i,j,k$ distinct.  If
  $\alpha\in[g_i,g_j]\cap[g_i,g_k]$, then
  \eqref{eq:geodesic-affine-line-constraint} gives, for every $x$,
  \begin{equation}
    \alpha(x)-x\in
    \bigl(e_i+\langle e_j-e_i\rangle\bigr)
    \cap
    \bigl(e_i+\langle e_k-e_i\rangle\bigr).
  \end{equation}
  This intersection is the single point $\{e_i\}$.  To see this, write
  \begin{equation}
    (1-t)e_i+te_j=(1-s)e_i+se_k,\qquad s,t\in\mathbb Z_n.
  \end{equation}
  Since the vectors $e_1,\ldots,e_{q-1},e_q=0$ are the vertices of the
  standard simplex, comparing the coordinates indexed by
  $\{i,j,k\}\setminus\{q\}$ gives $s=t=0$.  Thus
  $\alpha(x)=x+e_i=g_i(x)$ for every $x$, and hence
  $\alpha=g_i$.  This proves
  $[g_i,g_j]\cap[g_i,g_k]=\{g_i\}$, and the same argument gives the
  stated singleton $\{g_m\}$ for any common terminal $m$.

  Finally suppose that $\{i,j\}\cap\{k,\ell\}=\varnothing$.  If
  $\alpha$ belonged to both intervals, then for every $x$ the
  displacement $\alpha(x)-x$ would lie in
  \begin{equation}
    \bigl(e_i+\langle e_j-e_i\rangle\bigr)
    \cap
    \bigl(e_k+\langle e_\ell-e_k\rangle\bigr).
  \end{equation}
  This intersection is empty.  Indeed an equality
  \begin{equation}
    (1-t)e_i+te_j=(1-s)e_k+se_\ell
  \end{equation}
  would give two affine combinations of the simplex vertices with total
  coefficient $1$.  Comparing the coordinates of
  $e_1,\ldots,e_{q-1}$, and then using the total coefficient, gives
  uniqueness of these affine coordinates.  This is impossible because
  the two unordered pairs of vertices are disjoint.  Hence the two
  intervals are disjoint.
\end{proof}

\end{document}
